% Options macros, minted by experiments/spxw_options_tex.py. Each is a
% \newcommand, so the file is read exactly once however the document is
% assembled; the guard makes a second \input elsewhere a no-op.
\ifdefined\optNContracts\else% GENERATED FILE -- do not edit by hand.
% Written by experiments/spxw_options_tex.py from results/spxw_pnl/*.csv.
% Options-section macros: the delta-hedged 0DTE leg book, its cost floor,
% regime/exit rules, and the time-till-close forecast structure.
% MODE: --ft. Every macro in this file is F_t-measurable: the
% dh-family macros quote the strictly-at-entry ledger, and the
% toclose/encompassing macros quote the at-entry remaining
% forecast F_rem (experiments/ft_remaining.py). The *_rem sums
% (forecasts issued during the trade window) survive only as
% labelled ex-post decomposition rows in the CSVs ('pRem ...' /
% unprefixed 'paper ...' rows) and feed no macro.
\newcommand{\optAlwaysShortHit}{66\%}
\newcommand{\optAlwaysShortSh}{4.26}
\newcommand{\optBareAzeroHit}{49\%}
\newcommand{\optBareAzeroLong}{34\%}
\newcommand{\optBareAzeroShMid}{$-$0.1}
\newcommand{\optBareAzeroShX}{$-$3.4}
\newcommand{\optBareAzeroTraded}{60\%}
\newcommand{\optBareBlkHit}{51\%}
\newcommand{\optBareBlkLong}{33\%}
\newcommand{\optBareBlkShMid}{0.7}
\newcommand{\optBareBlkShX}{$-$2.5}
\newcommand{\optBareBlkTraded}{59\%}
\newcommand{\optBareCallAzeroShMid}{$-$0.6}
\newcommand{\optBareCallAzeroShX}{$-$3.9}
\newcommand{\optBareCallBlkShMid}{0.2}
\newcommand{\optBareCallBlkShX}{$-$2.9}
\newcommand{\optBareCtrlHit}{66\%}
\newcommand{\optBareCtrlLong}{0\%}
\newcommand{\optBareCtrlShMid}{4.3}
\newcommand{\optBareCtrlShX}{---}
\newcommand{\optBareCtrlTraded}{100\%}
\newcommand{\optBareDailyAzeroHit}{47\%}
\newcommand{\optBareDailyAzeroLong}{39\%}
\newcommand{\optBareDailyAzeroShMid}{$-$1.3}
\newcommand{\optBareDailyAzeroShX}{$-$4.5}
\newcommand{\optBareDailyAzeroTraded}{62\%}
\newcommand{\optBareDailyBlkHit}{50\%}
\newcommand{\optBareDailyBlkLong}{35\%}
\newcommand{\optBareDailyBlkShMid}{$-$0.4}
\newcommand{\optBareDailyBlkShX}{$-$3.4}
\newcommand{\optBareDailyBlkTraded}{60\%}
\newcommand{\optBareDailyCallAzeroShMid}{$-$1.6}
\newcommand{\optBareDailyCallAzeroShX}{$-$4.8}
\newcommand{\optBareDailyCallBlkShMid}{$-$0.7}
\newcommand{\optBareDailyCallBlkShX}{$-$3.6}
\newcommand{\optBareDailyCtrlHit}{68\%}
\newcommand{\optBareDailyCtrlLong}{0\%}
\newcommand{\optBareDailyCtrlShMid}{4.7}
\newcommand{\optBareDailyCtrlShX}{---}
\newcommand{\optBareDailyCtrlTraded}{100\%}
\newcommand{\optBareDailyPutAzeroShMid}{$-$1.0}
\newcommand{\optBareDailyPutAzeroShX}{$-$4.2}
\newcommand{\optBareDailyPutBlkShMid}{$-$0.3}
\newcommand{\optBareDailyPutBlkShX}{$-$3.2}
\newcommand{\optBarePutAzeroShMid}{0.2}
\newcommand{\optBarePutAzeroShX}{$-$2.9}
\newcommand{\optBarePutBlkShMid}{1.0}
\newcommand{\optBarePutBlkShX}{$-$2.1}
\newcommand{\optBothShMid}{0.7}
\newcommand{\optBothShX}{$-$2.5}
\newcommand{\optBumpsCtrlShMid}{3.7}
\newcommand{\optBumpsCtrlShX}{1.1}
\newcommand{\optBumpsXsMidHi}{0.9}
\newcommand{\optBumpsXsMidLo}{$-$0.6}
\newcommand{\optBumpsXsXHi}{$-$1.2}
\newcommand{\optBumpsXsXLo}{$-$6.6}
\newcommand{\optCallShMid}{0.2}
\newcommand{\optCallShX}{$-$2.9}
\newcommand{\optCostDeepHi}{0.14}
\newcommand{\optCostDeepLo}{0.08}
\newcommand{\optCostElevenLo}{0.018}
\newcommand{\optCostElevenMid}{0.024}
\newcommand{\optCostFifteenLo}{0.032}
\newcommand{\optCostFifteenMid}{0.049}
\newcommand{\optDailyBothShMid}{$-$0.4}
\newcommand{\optDailyBothShX}{$-$3.4}
\newcommand{\optDailyCallShMid}{$-$0.7}
\newcommand{\optDailyCallShX}{$-$3.6}
\newcommand{\optDailyNDays}{337}
\newcommand{\optDailyPutShMid}{$-$0.3}
\newcommand{\optDailyPutShX}{$-$3.2}
\newcommand{\optEncBlkDM}{$-$4.1}
\newcommand{\optEncBlkQLIKE}{0.1503}
\newcommand{\optEncBlkRawDM}{$-$3.3}
\newcommand{\optEncBlkRawQLIKE}{0.1653}
\newcommand{\optEncJointDM}{$-$5.1}
\newcommand{\optEncJointQLIKE}{0.1381}
\newcommand{\optEncMfivDM}{---}
\newcommand{\optEncMfivQLIKE}{0.2374}
\newcommand{\optEncNDays}{545}
\newcommand{\optExitCrossEarly}{67\%}
\newcommand{\optExitCrossShMid}{1.66}
\newcommand{\optExitCrossShX}{$-$2.92}
\newcommand{\optExitHoldShMid}{0.67}
\newcommand{\optExitHoldShX}{$-$2.47}
\newcommand{\optExitKFourShMid}{0.59}
\newcommand{\optExitKFourShX}{$-$6.91}
\newcommand{\optExitKOneShMid}{2.08}
\newcommand{\optExitKOneShX}{$-$13.14}
\newcommand{\optFinalAllCrossHitX}{42\%}
\newcommand{\optFinalAllCrossNDays}{565}
\newcommand{\optFinalAllCrossShMid}{1.66}
\newcommand{\optFinalAllCrossShX}{$-$2.92}
\newcommand{\optFinalAllHoldHitX}{46\%}
\newcommand{\optFinalAllHoldNDays}{565}
\newcommand{\optFinalAllHoldShMid}{0.67}
\newcommand{\optFinalAllHoldShX}{$-$2.47}
\newcommand{\optFinalDailyShMid}{0.33}
\newcommand{\optFinalDailyShX}{$-$4.04}
\newcommand{\optFinalExCrossHitX}{40\%}
\newcommand{\optFinalExCrossNDays}{549}
\newcommand{\optFinalExCrossShMid}{1.55}
\newcommand{\optFinalExCrossShX}{$-$3.07}
\newcommand{\optFinalExHoldHitX}{44\%}
\newcommand{\optFinalExHoldNDays}{549}
\newcommand{\optFinalExHoldShMid}{0.55}
\newcommand{\optFinalExHoldShX}{$-$2.60}
\newcommand{\optFinalFiveShMid}{1.71}
\newcommand{\optFinalFiveShX}{$-$3.32}
\newcommand{\optFinalFiveTraded}{78\%}
\newcommand{\optFinalSwapTFiveHund}{+2.06}
\newcommand{\optFinalSwapTTen}{+2.05}
\newcommand{\optFinalSwapTTwenty}{+2.19}
\newcommand{\optFomcDays}{16}
\newcommand{\optFomcFracShort}{98\%}
\newcommand{\optFomcShMid}{4.81}
\newcommand{\optFomcShX}{2.33}
\newcommand{\optFomcUnknownDays}{57}
\newcommand{\optForecastMode}{the one-step forecast standing at $t$, mapped to the remaining horizon}
\newcommand{\optLateLong}{23\%}
\newcommand{\optLateShHi}{1.7}
\newcommand{\optLateVrp}{0.06}
\newcommand{\optMornCtrlHi}{4.8}
\newcommand{\optMornCtrlLo}{3.6}
\newcommand{\optMornShHi}{1.1}
\newcommand{\optMornShLo}{0.1}
\newcommand{\optNContracts}{60{,}532}
\newcommand{\optNDays}{565}
\newcommand{\optNoFomcShMid}{0.84}
\newcommand{\optNoFomcShX}{$-$2.25}
\newcommand{\optPutShMid}{1.0}
\newcommand{\optPutShX}{$-$2.1}
\newcommand{\optRealizedKFour}{0\%}
\newcommand{\optRealizedKNine}{71\%}
\newcommand{\optRealizedKOne}{49\%}
\newcommand{\optSmileShortSh}{8.6}
\newcommand{\optStripShortSh}{6.7}
\newcommand{\optThetaFiveHundAzeroShMid}{0.1}
\newcommand{\optThetaFiveHundAzeroShX}{$-$3.6}
\newcommand{\optThetaFiveHundHit}{51\%}
\newcommand{\optThetaFiveHundLong}{43\%}
\newcommand{\optThetaFiveHundShMid}{0.9}
\newcommand{\optThetaFiveHundShX}{$-$2.6}
\newcommand{\optThetaFiveHundSwapT}{+2.83}
\newcommand{\optThetaFiveHundTraded}{79\%}
\newcommand{\optThetaTenAzeroShMid}{$-$0.1}
\newcommand{\optThetaTenAzeroShX}{$-$3.4}
\newcommand{\optThetaTenHit}{51\%}
\newcommand{\optThetaTenLong}{33\%}
\newcommand{\optThetaTenShMid}{0.7}
\newcommand{\optThetaTenShX}{$-$2.5}
\newcommand{\optThetaTenSwapT}{+3.00}
\newcommand{\optThetaTenTraded}{59\%}
\newcommand{\optThetaTwentyAzeroShMid}{$-$0.1}
\newcommand{\optThetaTwentyAzeroShX}{$-$2.5}
\newcommand{\optThetaTwentyHit}{52\%}
\newcommand{\optThetaTwentyLong}{15\%}
\newcommand{\optThetaTwentyShMid}{0.5}
\newcommand{\optThetaTwentyShX}{$-$2.1}
\newcommand{\optThetaTwentySwapT}{+3.33}
\newcommand{\optThetaTwentyTraded}{30\%}
\newcommand{\optThetaZeroAzeroShMid}{0.2}
\newcommand{\optThetaZeroAzeroShX}{$-$3.9}
\newcommand{\optThetaZeroHit}{50\%}
\newcommand{\optThetaZeroLong}{54\%}
\newcommand{\optThetaZeroShMid}{1.0}
\newcommand{\optThetaZeroShX}{$-$2.8}
\newcommand{\optThetaZeroSwapT}{+2.66}
\newcommand{\optThetaZeroTraded}{100\%}
\newcommand{\optTocloseAzeroQLIKE}{0.1635}
\newcommand{\optTocloseBlkQLIKE}{0.1653}
\newcommand{\optTocloseHessHit}{59\%}
\newcommand{\optTocloseHessSh}{2.1}
\newcommand{\optTocloseIncrSh}{$-$0.3}
\newcommand{\optTocloseNDays}{671}
\newcommand{\optTtcABlkDM}{---}
\newcommand{\optTtcABlkNested}{---}
\newcommand{\optTtcABlkQLIKE}{0.0634}
\newcommand{\optTtcAaZeroDM}{+2.4}
\newcommand{\optTtcAaZeroNested}{---}
\newcommand{\optTtcAaZeroQLIKE}{0.0693}
\newcommand{\optTtcBFtBlkDM}{+8.7}
\newcommand{\optTtcBFtBlkNested}{---}
\newcommand{\optTtcBFtBlkQLIKE}{0.2137}
\newcommand{\optTtcBFtjointDM}{+9.1}
\newcommand{\optTtcBFtjointNested}{$-$3.6}
\newcommand{\optTtcBFtjointQLIKE}{0.1912}
\newcommand{\optTtcBdirectBlkDM}{+1.0}
\newcommand{\optTtcBdirectBlkNested}{---}
\newcommand{\optTtcBdirectBlkQLIKE}{0.0642}
\newcommand{\optTtcBdirectIvDM}{+15.0}
\newcommand{\optTtcBdirectIvNested}{---}
\newcommand{\optTtcBdirectIvQLIKE}{0.3108}
\newcommand{\optTtcBjointDM}{$-$0.7}
\newcommand{\optTtcBjointNested}{$-$3.6}
\newcommand{\optTtcBjointQLIKE}{0.0627}
\newcommand{\optTtcBjointrealizedDM}{$-$2.8}
\newcommand{\optTtcBjointrealizedNested}{$-$3.1}
\newcommand{\optTtcBjointrealizedQLIKE}{0.0577}
\newcommand{\optTtcDailyABlkDM}{---}
\newcommand{\optTtcDailyABlkNested}{---}
\newcommand{\optTtcDailyABlkQLIKE}{0.0626}
\newcommand{\optTtcDailyAaZeroDM}{+2.1}
\newcommand{\optTtcDailyAaZeroNested}{---}
\newcommand{\optTtcDailyAaZeroQLIKE}{0.0698}
\newcommand{\optTtcDailyBFtBlkDM}{+6.3}
\newcommand{\optTtcDailyBFtBlkNested}{---}
\newcommand{\optTtcDailyBFtBlkQLIKE}{0.2295}
\newcommand{\optTtcDailyBFtjointDM}{+6.4}
\newcommand{\optTtcDailyBFtjointNested}{$-$4.2}
\newcommand{\optTtcDailyBFtjointQLIKE}{0.1900}
\newcommand{\optTtcDailyBdirectBlkDM}{+1.9}
\newcommand{\optTtcDailyBdirectBlkNested}{---}
\newcommand{\optTtcDailyBdirectBlkQLIKE}{0.0631}
\newcommand{\optTtcDailyBdirectIvDM}{+11.3}
\newcommand{\optTtcDailyBdirectIvNested}{---}
\newcommand{\optTtcDailyBdirectIvQLIKE}{0.2126}
\newcommand{\optTtcDailyBjointDM}{$-$2.7}
\newcommand{\optTtcDailyBjointNested}{$-$4.3}
\newcommand{\optTtcDailyBjointQLIKE}{0.0606}
\newcommand{\optTtcDailyBjointrealizedDM}{$-$3.8}
\newcommand{\optTtcDailyBjointrealizedNested}{$-$3.8}
\newcommand{\optTtcDailyBjointrealizedQLIKE}{0.0519}
\newcommand{\optTtcNDays}{800}
\fi

\subsection{Reading the Forecast Against the Option Market}
\label{sec:options_methods}

% Methods only: instrument, measurement, signal, execution, controls.
% Every number here is a design constant or a measurement on the quote
% surface; no result is reported in this subsection.

The same-day index option market posts, at every half hour of every
expiration session, a price for the object this paper forecasts: the
variance still to come before the close. The forecast can therefore be
read against a quoted price as well as against a squared error. What
follows fixes the instrument, the two measurements that have to be made
commensurate, the decision rule, the execution convention, and the
controls; every constant it names was set before the book was scored.

\paragraph{Instrument and panel.} S\&P~500 index (SPX) options with zero
days to expiry, cash settled at the 16:00~ET close, so a position carried
to settlement is closed at intrinsic with no share unwinding. The panel
is every expiration-day chain from 2020 through 2025 at 30-minute stamps
within regular hours (bid, ask, mid, underlying), joined on the ET wall
clock to the paper's per-bar forecasts --- the OLS--HAR incumbent of
Section~\ref{sec:dp_benchmark} and the two-block ridge of
Section~\ref{sec:block_tikhonov} --- and to realized per-bar variance
from the panel of Section~\ref{sec:data_prep}. A contract enters when its
Black--Scholes $|\Delta|$ lies in the near-the-money band $[0.30, 0.70]$
and both sides are quoted ($\mathrm{bid} > 0$): \optNContracts{} contract
entries over \optNDays{} expiration days. Calls and puts are treated
separately rather than combined into straddles --- the sample is larger,
mispricing is visible leg by leg instead of netted within a strike, and
the two legs' agreement is itself a check that the signal is about
variance and not direction. Entries begin at 10:00, because the vendor
supplies neither underlying price nor greeks at the 09:30 stamp, and half
sessions --- identified by disagreement between the 16:00 clock and the
vendor's own hours-to-expiration --- are dropped whole.

\paragraph{Time-till-close implied variance.} At each stamp the implied
volatility is solved from the mid under Black--Scholes with $r = 0$ and
$\tau$ the trading time remaining to the settle,
\begin{equation}
    V^{\mathrm{BS}}(S, K, \tau, \sigma)
    = \phi \bigl[\, S\,\Phi(\phi d_1) - K\,\Phi(\phi d_2) \,\bigr],
    \quad
    d_1 = \frac{\log (S/K) + \tfrac{1}{2}\sigma^2 \tau}{\sigma\sqrt{\tau}},
    \quad
    d_2 = d_1 - \sigma\sqrt{\tau},
    \label{eq:opt_bs}
\end{equation}
with $\phi = +1$ for a call and $-1$ for a put and
$\tau = (\text{hours to } 16{:}00)/(252 \times 6.5)$ on a DST-aware
clock; $\mathrm{IV}_{t,T}$ is the root of Equation~\eqref{eq:opt_bs} at
the quoted mid, and the vendor's own volatility column is never read.
Black--Scholes is sufficient here, and the reason is worth stating
because the choice would otherwise read as an assumption. The position
settles at intrinsic, a model-free payoff: realized profit and loss is
fixed by the settlement price and the hedge path, not by a pricing model.
Black--Scholes enters only as the quoting map that renders prices
comparable across strikes and stamps as a time-till-close volatility, and
as the hedge ratio, a second-order term; any smile-consistent model gives
the identical intrinsic profit and loss and differs only in the delta it
prescribes. Where a strike-independent measurement is wanted the
model-free log strip is used instead --- out-of-the-money puts below the
forward, calls above, the at-the-money strike averaged across the two
sides, trapezoidal strike weights,
\begin{equation}
    C_t \;=\; 2 \sum_i \frac{\Delta K_i}{K_i^{2}}\, Q_t(K_i)
    \;-\; \Bigl( \frac{F_t}{K_0} - 1 \Bigr)^{2},
    \label{eq:opt_mfiv}
\end{equation}
the integrated implied variance to the close, written MFIV below. It is a
strip mark rather than a tradeable single contract, and it is used for
the encompassing and information-set tests only.

\paragraph{Time-till-close variance forecast.} The model's counterpart to
$C_t$ is the sum of its per-bar forecasts over the remaining bars of the
session, rescaled by a causally estimated multiplicative factor,
\begin{equation}
    \widehat{RV}_{t,T} \;=\; c_{h(t)} \sum_{t < t_k \le T} \hat v_{t_k},
    \qquad
    c_{h} \;=\; \operatorname*{mean}_{d' < d}
      \frac{RV_{d',h}}{\sum_k \hat v_{d',k}},
    \label{eq:opt_ttc}
\end{equation}
where $h(t)$ is the entry hour, the mean of realized-to-forecast ratios
is the QLIKE-optimal multiplicative rescale, and the average runs over
strictly earlier expiration days with a minimum of $63$ before any value
is issued, applied shifted one day. The result is converted to an
annualized volatility on the option's own horizon,
$\hat\sigma_{t,T} = ( \widehat{RV}_{t,T} / \tau_t )^{1/2}$, the unit
$\mathrm{IV}_{t,T}$ already carries. One caveat is stated here and not
repeated: the per-bar forecast for a later bar is formed \emph{at} that
bar, so the sum is not $\mathcal{F}_t$-measurable and quietly carries
information the trader standing at $t$ does not have. It is the
construction used elsewhere in this paper and a generous one. The
strictly $\mathcal{F}_t$ alternatives --- the one-step forecast standing
at $t$, and a direct regression of remaining variance on information
available at $t$ --- are reported alongside it in the forecast-structure
comparison over \optTtcNDays{} days; the trading book is a demonstration
run under the sum.

\paragraph{Signal, dead zone, and the cost floor.} The signal is the
price-form variance risk premium, the log ratio of the two commensurate
volatilities, and the rule is a dead zone around zero:
\begin{equation}
    s_t \;=\; \log \frac{\mathrm{IV}_{t,T}}
                        {\sqrt{\widehat{RV}_{t,T} / \tau_t}} ,
    \qquad
    q_t \;=\; -\,\mathbf{1}\{ s_t > \theta \}
              \;+\; \mathbf{1}\{ s_t < -\theta \},
    \label{eq:opt_signal}
\end{equation}
so the book sells the contract and hedges it when the option is expensive
relative to the forecast by more than $\theta$, buys and hedges when it
is cheap by more than $\theta$, and stands flat in between. The dead zone
is read from a grid fixed in advance, $\theta \in \{0, 0.05, 0.10,
0.20\}$, every level reported; no value is selected on realized profit
and loss. Its floor is a measurement, not a preference: crossing the
spread moves the quoted volatility, and in the same log-volatility unit
as $s_t$ the round trip costs
\begin{equation}
    \theta_{\min}
    \;=\; \log \frac{\mathrm{IV}(\mathrm{ask})}{\mathrm{IV}(\mathrm{bid})},
    \label{eq:opt_costfloor}
\end{equation}
both volatilities solved from Equation~\eqref{eq:opt_bs} at the same
stamp and strike. The median floor is \optCostElevenLo{} at 11:00~ET for
$|\Delta| \in [0.30, 0.50)$ and \optCostElevenMid{} for $[0.50, 0.70)$,
rising to \optCostFifteenLo{} and \optCostFifteenMid{} by 15:00~ET; for
$|\Delta| > 0.70$ it is \optCostDeepLo{}--\optCostDeepHi{}, which no
signal in this study approaches, and that measurement rather than a taste
for at-the-money exposure is why the band stops at $0.70$. The absolute
half-spread is pinned at the minimum tick across the session, so the
intraday rise in $\theta_{\min}$ is the premium thinning toward the close
against a fixed tick, not quotes widening. A dead zone below the measured
floor is one the market has already closed.

\paragraph{Delta hedge and settlement.} Each entered contract is hedged
with the Black--Scholes delta evaluated at that contract's \emph{entry}
implied volatility and rebalanced at every subsequent 30-minute stamp on
the chain's own underlying. With $\delta_k$ that delta at stamp $t_k$,
$S_k$ the underlying, and $V_{t_0}$ the entry mid, the normalized profit
and loss of a unit position is
\begin{equation}
    \Pi \;=\; \frac{q}{V_{t_0}}
      \Bigl[\, \bigl(\phi (S_T - K)\bigr)^{+} - V_{t_0}
      \;-\; \sum_{k} \delta_k \bigl( S_{k+1} - S_k \bigr) \Bigr],
    \label{eq:opt_pnl}
\end{equation}
the option leg settled at intrinsic and everything normalized by the
entry mid so contracts of different premium are comparable within a day.
Two execution conventions run side by side: a mid-quote variant, and a
crossed variant charging the option at the bid when sold and at the ask
when bought, with $0.5$~bp of traded notional on the underlying per
rebalance --- a documented assumption, since no trade sizes exist on
disk. Two exit rules likewise run side by side: hold to settlement, and
exit at the first later stamp at which the re-measured signal crosses
back through zero, which pays the option spread again in the crossed
variant. Neither is selected on outcome.

\paragraph{Aggregation and controls.} Contracts within a session are not
independent draws, so the unit of account is the day: one
profit-and-loss number per expiration day, the mean over that day's
traded contracts. The annualized Sharpe ratio is the mean of those daily
numbers over their standard deviation, scaled by $\sqrt{252}$; the hit
rate is the fraction of positive days. The book is cut per leg, per era
(the whole sample against the daily-0DTE era beginning 16~May~2022, from
which every weekday carries an expiration), and per entry hour. Two
controls run on the identical contract set. The first is unconditional:
always short, and always long, the same contracts under the same hedge
and settlement conventions --- the variance risk premium with no forecast
in it, which is what any conditional book must be measured against. The
second is the \emph{swap test}: the identical machinery with the OLS--HAR
incumbent's forecast substituted for the two-block ridge in
Equation~\eqref{eq:opt_ttc}, compared by a paired daily difference
$t$-statistic on the same days. Because the two arms differ only in the
forecast entering the signal, the swap test is the test of whether the
exogenous block, rather than the HAR family it is attached to, is doing
the work. As robustness the book is tabulated by causal VIX quintile, VIX
term slope, day of week, macro-release day, and FOMC-announcement day;
the release flags end 2024-01-09, and later days are \emph{unknown}
rather than ``no announcement.''

\paragraph{Encompassing and information-set tests.} The book asks whether
the forecast is worth money; the regressions ask the prior question of
whether it carries information the market price does not. Two are run: at
10:00~ET, an expanding log-log regression of the variance realized to the
close on the strip alone, on the model alone, and on both jointly; and at
every stamp, a direct expanding regression of remaining variance on the
information available at that stamp, built up one regressor at a time,
\begin{align}
    \log RV_{10{:}00,T}
    &= a + b \log C_{10{:}00} + c \log \widehat{RV}_{10{:}00,T}
    + \varepsilon,
    \label{eq:opt_enc} \\
    \log RV_{t,T}
    &= a_h + b_h \log \widehat{RV}^{\mathrm{sum}}_{t,T}
    + c_h \log C_t + d_h \log RV_{\mathrm{open},t} + \varepsilon_t,
    \label{eq:opt_direct}
\end{align}
where $RV_{\mathrm{open},t}$ is the variance already realized in the
session. Equation~\eqref{eq:opt_direct} is fit two ways --- separately
per entry hour, and once pooled with hour dummies --- so that the gain
from hour-specific coefficients is separated from the gain from the
regressors. Strictly $\mathcal{F}_t$ versions of every variant replace
the remaining sum with the one-step forecast standing at $t$, which
together with $C_t$ and $RV_{\mathrm{open},t}$ makes the whole right-hand
side measurable at the decision time. Every fit is causal on the
convention of Equation~\eqref{eq:opt_ttc} --- expanding over strictly
earlier days, minimum $63$, applied one day forward --- back-transformed
by the expanding multiplicative factor computed in the same grouping as
the fit it corrects, scored by QLIKE \citep{Patton2011} on remaining
variance, and compared by a Diebold--Mariano statistic
\citep{DieboldMariano1995,HarveyLeybourneNewbold1997} on paired daily
loss differences.

\paragraph{What is not modelled.} Two things lie outside what this data
can support, and neither is approximated. The first is capacity: no trade
sizes exist on disk, so Equation~\eqref{eq:opt_costfloor} and the crossed
variant are a \emph{quote-cost} study at the posted top of book, and
depth, queue position, and market impact are not measurable here. The
second is cross-expiry structure: only expiration-day chains are on disk,
and within a single expiry the vega-to-gamma ratio is $S^2 \sigma \tau$
at every strike, so no position in that expiry is long gamma and flat
vega and the gamma exposure cannot be separated from the marks the smile
assigns it.
